Volumetric media such as \glsentrylong{PC}s and meshes are the basis of emerging virtual and augmented reality experiences, but practical delivery relies on lossy compression that can visibly degrade quality. This monograph develops \glsentrylong{FR} quality assessment methods, which compare a distorted asset against its pristine reference, to estimate perceived quality with minimal computational cost. The goal is to give practitioners tools that reflect human opinion while remaining fast enough for real-world pipelines.

We explore three complementary paradigms. Projection-based approaches render volumetric content into orthographic cubemaps and apply mature deep \glsentrylong{IQA} metrics to the resulting images. Descriptor-based approaches analyze local geometry and color through \glsentrylong{PCA}-driven statistics. Hybrid approaches fuse both views to capture appearance and structure simultaneously.

Four main results stem from this exploration. (i) A projection-based \glsentrylong{PC} metric that attains high correlation with subjective scores using deep \glsentrylong{IQA} on post-processed cubemaps. (ii) PointPCA+RS, a Rust reimplementation of PointPCA+ focused on high computational performance, delivering up to $29\times$ faster runtime and $33\times$ lower memory usage while matching accuracy. (iii) PointPCA++, a hybrid \glsentrylong{PC} metric that joins projection features with optimized \glsentrylong{PCA}-based descriptors to improve robustness across distortion types. (iv) \glsentrylong{MAPPE}, a mesh metric that unifies a custom Rust renderer, a \glsentrylong{TAPS} strategy, projection scoring, and an efficient PointPCA\textsuperscript{3} descriptor module for dense mesh-derived \glsentrylong{PC}s.

Across public datasets, the proposed metrics achieve \glsentrylong{SROCC} above 0.90 and, in the mesh case, outperform prior hybrid methods while being over two orders of magnitude faster than the closest baseline. These results indicate that perceptually aligned, \glsentrylong{FR} volumetric quality assessment can be both accurate and practical for immersive media workflows.
